\section{Numerical Simulations and Results}
\label{sec:numerics}

In this section, we present our lattice QCD results. We start with the lattice setup, followed by results
for quasi TMDWFs with two-point correlations. 
The Wilson loop results are discussed in subsection III.C. Subsection III.D studies the operator mixing effects. Our main result on CS kernel is presented in subsection E. The final subsection E includes some overall discussions. 

\subsection{Lattice setup }

Our numerical simulations use $N_f=2+1+1$ valence clover fermions on a  highly improved staggered quark (HISQ) sea~\cite{Follana:2006rc} and a 1-loop Symanzik improved gauge action~\cite{Symanzik:1983dc}, generated by the MILC collaboration~\cite{MILC:2012znn} using periodic boundary conditions. In the calculations, we use a single ensemble with the lattice spacing $a\simeq0.12$~fm and the volume $n_s^3\times n_t=48^3\times64$ at physical sea-quark masses. In order to increase the signal-to-noise ratio, we tune the light-valence quark masses to the strange-quark one, namely $m_{\pi}^{\mathrm{sea}}=130$~MeV and  $m_{\pi}^{\mathrm{val}}=670$~MeV, which
could generate some non-unitarity effects. On the other hand, the Collins-Soper kernel only depends weakly on quark mass, and we may consider valence quarks are strange-like, namely the hadrons involved are kaons.  

To further improve the statistical signals, we adopt hypercubic (HYP) smeared fat links~\cite{Hasenfratz:2001hp} for the gauge ensembles. To access the large momentum limit for the CS kernel, we employ three different hadron momenta as $P^z=2\pi/n_s\times\{8,10,12\}=\{1.72,~2.15,~2.58\}$~GeV  corresponding  to the boost factor $\gamma=\{2.57,~3.21,~3.85\}$. 

%%%%%%%%%%%%%%%%%%%%%%%%%%

\subsection{Quasi TMDWFs From Two-Point Correlators}
\label{sec:quasi-WF_from_2pt}
In order to calculate the quasi TMDWFs  defined in Eq.(\ref{eq:quasiWFinmomentumspace}), we generate Coulomb-gauge wall-source propagators, 
\begin{align}
S_{w}\left(x, t, t^{\prime} ; \vec{p}\right)=\sum_{\vec{y}} S\left(t, \vec{x} ; t^{\prime}, \vec{y}\right) e^{i \vec{p} \cdot(\vec{y}-\vec{x})}, \label{eq:wallsource}
\end{align}
where $(t^{\prime},\vec{y})$ and $(t,\vec{x})$ denote the space-time positions of source and sink. Then one  can construct the two-point function (2pt) related to the quasi TMDWFs  in Eq.~(\ref{eq:quasiWFinmomentumspace}):
\begin{align}
C_2^{\pm}&	\left(z,b_{\perp},P^z; p^z,L,t \right)=\frac{1}{n_s^3}\sum_{\vec{x}}\mathrm{tr} e^{i\vec{P}\cdot\vec{x}}\left\langle S_w^{\dagger}\left(\vec{x}_1,t,0;-\vec{p}\right)\right. \nonumber\\
&\;\;\;\;\;\;\;\;\;\; \;\;\;\;\;\;\;\;\;\;    \times \left. \Gamma U_{\sqsupset, \pm L}\left(\vec{x}_1,\vec{x}_2\right) S_w\left(\vec{x}_2,t,0;\vec{p} \right)  \right\rangle
\label{eq:nonlocal2pt}
\end{align}
with $\vec{x}_1=\vec{x}+z\hat{n}_z/2+b_{\perp}\hat{n}_{\perp}$ and  $\vec{x}_2=\vec{x}-z \hat{n}_{z}/2$. The quark momentum $\vec{p}=(\vec{0}_{\perp},p^z)$ is along the $z$-direction, and each of two quarks carries half of the hadron momentum. Thereby the hadron momentum satisfies  $\vec{P}=2\vec{p}$. The anti-quark propagator can be obtained from Eq.(\ref{eq:1loophardkernel}) by applying $\gamma_5$-hermiticity $S_w(x,y)=\gamma_5S_w(y,x)^{\dagger}\gamma_5$. As mentioned above, the Dirac structures are chosen as $\Gamma=\gamma^z\gamma_5$ and $\gamma^t\gamma_5$ that can be projected onto the leading twist light-cone contributions in the large $P^z$ limit.
 
By generating the wall-source propagators with quark momenta $p^z=\pm\{4,5,6\}\times2\pi/(n_sa)$, and three segments of gauge links following Eq.(\ref{eq:stapleshapedUlink}), one  can construct  the two-point correlation functions on the  lattice. With the help of reduction formulas, the $C_2^{\pm}$ can be parametrized as
 \begin{align}
C_2^{\pm}&\left(z,b_{\perp},P^z; p^z,L,t \right)=\frac{A_w(p^z)A_p}{2E}\tilde{\Phi}^{\pm0}\left(z,b_{\perp},P^z,L\right) \nonumber\\
& \times e^{-Et}\left[1+c_0\left(z,b_{\perp},P^z,L\right)e^{-\Delta Et} \right], \label{eq:C2parametrization}
 \end{align}
where $A_w(p^z)$ is the matrix element of the pseudoscalar meson interpolating field with Coulomb gauge fixed wall source and $A_p$ is the one for a point source (sink). These terms as well as the factor $E^{-1}=1/\sqrt{m_{\pi}^2+\left(P^z\right)^2}$ are cancelled by the local  two-point function $C_2^{\pm}\left(0,0,P^z; p^z,0,t \right)$ at the same time slice. Thus the remaining  ground-state matrix element $\tilde{\Phi}^{\pm0}$ is normalized. The ratio of nonlocal and local two-point functions can be parametrized  as
\begin{align}
&R^{\pm}\left(z,b_{\perp},P^z,L,t \right)=\frac{C_2^{\pm}\left(z,b_{\perp},P^z,L,t \right)}{C_2\left(0,0,P^z,0,t \right)} \nonumber\\
=&\tilde{\Phi}^{\pm0}\left(z,b_{\perp},P^z,L\right)\left[1+c_0\left(z,b_{\perp},P^z,L\right)e^{-\Delta Et} \right].
\label{eq:two-state-fit}
\end{align}
 %}
where $C_2$ in the denominator is a local
correlator.
 
 %%%%%%%%%%%%%%%%%%%%%%%%%%%%%%%%%%
\begin{figure}
\centering
\includegraphics[width=0.62\textwidth]{t_dep_b2z0_r}
\caption{Comparison of  two-state fit  and  one-state fit to extract $\tilde{\Phi}^{\pm0}\left(z,b_{\perp},P^z,L\right)$. Taking $\{z,b_{\perp},P^z,L\}=\{0a, 2a, 24\pi/n_s, 6a\}$ as example, we can see that  the two-state fit works for $t\in[2a, 8a]$ while the one-state fit works for $t\in[5a, 8a]$. The fitted  results are consistent with each other, while  the one-state fit is more conservative. }
\label{fig:tdependenceofratio}
\end{figure}
%%%%%%%%%%%%%%%%%%%%%%%%%%%%%%%%%%
 
 
In the above parametrization,  the excited-state contributions are collected into the $c_0$ term, and $\Delta E$ denotes the mass gap between the ground and first excited state.  With the increase of Euclidean time, contributions from the excited state decay and the plateau obtained for $R^{\pm}\left(z,b_{\perp},P^z,L,t \right)$ at large times reflects the ground-state contribution $\tilde{\Phi}^{\pm0}$. We employ two methods to extract $\tilde{\Phi}^{\pm0}$, namely the two-state fit directly using Eq.~\eqref{eq:two-state-fit}, and the one-state fit by setting $c_0=0$. With a large enough Euclidean time, Fig.~\ref{fig:tdependenceofratio} exhibits a comparison using two methods for the case with small $\{z,b_{\perp}\}$. From this figure, one can see that the one-state fit result is consistent with two-state fit one but gives  a more conservative error estimate. The two-state fit works at $t\in[2a, 8a]$ and the one-state fit works at the plateau region $t\in[5a, 8a]$. While the excited state contamination would dominate for the two-state fit in a very high precision especially for the cases with large $\{z, b_\perp\}$. So with current accuracy, we adopt the more conservative results from the one-state fit in the following analysis. More details can be found in Appendix~\ref{appsec:C2}.



%%%%%%%%%%%%%%%%%%%%%%%%%%%%%%%%%%
\subsection{Wilson Loop Renormalization}
\label{subsec:numerical_wl}
%%%%%%%%%%%%%%%%%%%%%%%%%%%%%%%%%%

The unsubtracted quasi TMDWF matrix elements $\tilde{\Phi}^{\pm0}\left(z,b_{\perp},P^z,a,L\right)$ extracted from the joint fit of the  two-point function contain a factor $e^{-\delta \bar{m}(2L+b_{\perp})}$ from the linear divergence, the heavy quark effective potential factor $e^{-V(b_{\perp})L}$ and logarithmic divergences $Z_{O}$:
\begin{align}
\tilde{\Phi}^{\pm0}\left(z,b_{\perp},P^z,a,L\right) \propto e^{-\delta \bar{m}(2L+b_{\perp})}e^{-V(b_{\perp})L}Z_{O}.
\end{align} 
where $Z_O$ has logarithmic dependence on lattice spacing $a$. 

The linear divergence in $e^{-\delta \bar{m}(2L+b_{\perp})}$ comes from the self-energy of the Wilson line~\cite{Ji:2017oey,Ishikawa:2017faj,Green:2017xeu, Ji:2020brr}, where $\delta\bar m$ contains a term  proportional to $1/a$  and a  non-perturbative renormalon contribution $m_{0}$:
\begin{align}
\delta\bar m = \frac{m_{-1}(a)}{a} - m_0  \ . 
\end{align}
Note that the exponent of the linear divergence term is proportional to the total length of the Wilson link, e.g. $2L+b_{\perp}$ for the staple link. Due to this factor,  the numerical value for a Wilson loop dramatically decreases  for small $a$ and large $L$. 

The heavy quark effective potential term $e^{-V(b_{\perp})L}$ comes from interactions between the two Wilson lines along the $z$ direction in the staple link. The heavy quark effective potential $V(b_{\perp})$ is often used to determine the lattice spacing of an ensemble.  

The logarithmic divergence $Z_{O}$ comes from the vertices
involving the Wilson line and light quark. The logarithmic  divergence up to leading order, resumed by renormalization group equation, and matched to  $\overline{\rm MS}$ scheme is~\cite{Ji:1991pr, LatticePartonCollaborationLPC:2021xdx} 
\begin{align}\label{eq:ZO_RS_higher}
Z_{O}(1/a, \mu)=\left(\frac{\ln [1 /(a \Lambda_{\rm QCD}^{\rm latt})]}{\ln [\mu / \Lambda_{\rm QCD}^{\rm MS}]}\right)^{\frac{3 C_{F}}{b_0}},\nonumber\\
\end{align}
where $\Lambda_{\rm QCD}^{\rm latt}$  is different from that in $\overline{\rm MS}$ scheme. One can use both to effectively absorb higher order contributions~\cite{Lepage:1992xa}.

In this work,  the Wilson loop renormalization method~\cite{Chen:2016fxx,Zhang:2017bzy,Musch:2010ka,Green:2017xeu,Zhang:2017zfe}  is adopted, in which the Wilson loop $Z_E$ defined in Eq.(\ref{eq:WilsonLoop})  contains linear divergence and heavy quark effective potential:
\begin{align}
Z_E\left(2L, b_{\perp}, a\right) \propto e^{-\delta \bar{m}(4L+2b_{\perp})}e^{-V(b_{\perp})2L}.
\end{align}
According to Ref.~\cite{LatticePartonCollaborationLPC:2021xdx}, $\delta \bar{m}$ in the Wilson loop is the same as that in hadron matrix elements, and thus it is anticipated that the linear divergence is removed when dividing by $\sqrt{Z_E\left(2L, b_{\perp}, a\right)}$:
\begin{align}
\tilde{\Phi}^{\pm}(z,b_{\perp},P^z,a,L)=\frac{\tilde{\Phi}^{\pm0}\left(z,b_{\perp},P^z,a,L\right)}{\sqrt{Z_E\left(2L, b_{\perp}, a\right)}}.
\end{align}
As shown in Fig.~\ref{fig:WilsonLooprenormalization}, the subtracted quasi TMDWFs tend  to  be a constant when  $L\ge 0.4$ fm. We then use the subtracted quasi TMDWFs defined as
\begin{align}\label{eq:WLRQuasi}
\tilde{\Phi}^{\pm}(z,b_{\perp},P^z, a)=\lim_{L\to\infty}\frac{\tilde{\Phi}^{\pm0}\left(z,b_{\perp},P^z,a,L\right)}{\sqrt{Z_E\left(2L, b_{\perp}, a\right)}}.
\end{align}
However, it is anticipated that there is a residual logarithmic divergence $Z_{O}$:
\begin{align}
\frac{\tilde{\Phi}^{\pm0}\left(z,b_{\perp},P^z,a,L\right)}{\sqrt{Z_E\left(2L, b_{\perp}, a\right)}}
\propto \frac{e^{-\delta \bar{m}(2L+b_{\perp})}e^{-V(b_{\perp})L}Z_{O}}{e^{-\delta \bar{m}(2L+b_{\perp})}e^{-V(b_{\perp})L}}
= Z_{O}.
\end{align}
In the extraction of the CS kernel, a ratio of quasi TMDWFs is adopted and accordingly the residual logarithmic divergence $Z_{O}$ is canceled.



 %%%%%%%%%%%%%%%%%%%%%%%%%%%%%%%%%%
\begin{figure}
\centering
\includegraphics[width=0.62\textwidth]{l_dep_8_21_r}
\includegraphics[width=0.62\textwidth]{l_dep_8_21_i}
\caption{Results for the $L$-dependence of  unsubtracted and subtracted quasi TMDWFs:  real part (upper panel) and imaginary part (lower panel), as well as the square root of the Wilson loop. The case $\Gamma=\gamma^z\gamma_5$ and $\{P^z,b_{\perp}, z\}=\{16\pi/n_s,2a,2a\}$ is used for illustration. }
\label{fig:WilsonLooprenormalization}
\end{figure}
%%%%%%%%%%%%%%%%%%%%%%%%%%%%%%%%%%
 
\subsection{Higher-Twist Effects in Operators}
\label{sec:operator}


%%%%%%%%%%%%%%%%%%%%%%%%%%%%%%%%%%%%%%
\begin{figure}
\centering
\includegraphics[width=0.62\textwidth]{qwf_co_b3_r_p12}
\includegraphics[width=0.62\textwidth]{qwf_co_b3_i_p12}
\caption{$\lambda$-dependence of quasi-WF matrix elements with different Dirac structures. Here we take the case of $\{P^z,b_{\perp}\}=\{24\pi/n_s,3a\}$ as an example, the deviation between these two cases mainly comes from power corrections. Results for more sets of $\{P^z,b_{\perp}\}$ are shown in appendix~\ref{sec:app_operator_mixing}.}
    \label{fig:qwf_coordinate}
\end{figure}
%%%%%%%%%%%%%%%%%%%%%%%%%%%%%%%%%%%%%%


For a pseudoscalar meson on a Euclidean lattice, both $\Gamma=\gamma^t\gamma_5$ and $\gamma^z\gamma_5$ project onto the leading twist light-cone distribution amplitude, i.e. $\gamma^+\gamma_5$ in the arge $P^z$ limit. The differences  between them arises from power corrections in terms of $M^2/\left(P^z\right)^2$.

Fig.~\ref{fig:qwf_coordinate} shows the comparison of the $\lambda= z   P^z$-dependence of quasi TMDWFs with $\Gamma=\gamma^t\gamma_5$ and $\gamma^z\gamma_5$ at $P^z=24\pi/n_s\simeq2.58$~GeV. It can be seen from the plots that there are some differences between the two sets of results for the real part in  the small $\lambda$ region. The differences are expected to decrease with increasing $P^z$, and the correlators with $\gamma^t\gamma_5$ and $\gamma^z\gamma_5$ will gradually converge to the light-cone from opposite directions. Besides, in light-cone coordinate, $\gamma^t$ and $\gamma^z$ can be represented by $\gamma^+$ and $\gamma^-$,
\begin{align}
\gamma^t\gamma_5=\frac{1}{\sqrt{2}}(\gamma^++\gamma^-)\gamma_5,\nonumber\\
\gamma^z\gamma_5=\frac{1}{\sqrt{2}}(\gamma^+-\gamma^-)\gamma_5,
\end{align}
Due to the momentum along the light-cone, operators with $\gamma^-\gamma_5$ correspond to higher order terms of TMDWFs. Therefore, power corrections arising from finite $P^z$ are likely to be eliminated in the average of these two terms:
\begin{align}
\tilde{\Phi}^{\pm}=\frac{1}{2}\left[\tilde{\Phi}^{\pm}\left(\Gamma=\gamma^t\gamma_5\right) + \tilde{\Phi}^{\pm}\left(\Gamma=\gamma^z\gamma_5\right)\right].
\end{align}
For a quantitative analysis see the appendix of~\cite{LatticeParton:2020uhz}. The operator mixing effect reaches order 5\%~\cite{LatticeParton:2020uhz}, which is much smaller than the systematic uncertainties discussed in next subsection. 

According to our numerical simulations, the subtracted quasi TMDWFs in coordinate space $\tilde{\Phi}^{\pm}(z,b_{\perp},P^z)$ as a function of $\lambda=zP^z$ are complex, which is shown in Fig.~\ref{fig:qwf_co}. The examples are the real and imaginary part of $\tilde{\Psi}^{\pm}(x,b_{\perp},P^z)$ with $P^z=24\pi/n_s$, $b_{\perp}=2a$ and $4a$. 
To determine quasi TMDWFs in momentum space $\tilde{\Psi}^{\pm}(x,b_{\perp},P^z)$, we use a  Fourier transformation (FT)
\begin{align}
\tilde{\Psi}^{\pm}(x,b_{\perp},P^z)=\frac{1}{2\pi}\sum_{z_{\rm min}}^{z_{\rm max}}e^{ixzP^z}\tilde{\Phi}^{\pm}(z,b_{\perp},P^z).
\label{eq:FT}
\end{align}
Due to the imaginary part of $\tilde{\Phi}^{\pm}(z,b_{\perp},P^z)$,  $\tilde{\Psi}^{\pm}(x,b_{\perp},P^z)$ also has an imaginary part. We obtain the quasi TMDWFs in momentum space for both real and imaginary part of $\tilde{\Psi}^{\pm}(x,b_{\perp},P^z)$ shown in Fig.\ref{fig:qwf_momenta}, by taking $P^z=24\pi/n_s$ and $b_{\perp}=2a$ and $4a$ as examples. We truncate the FT at $z_{\rm min}$ and $z_{\rm max}$. The deviation of
  $\tilde{\Phi}^{\pm}(z_{\rm min},b_{\perp},P^z)$ and $\tilde{\Phi}^{\pm}(z_{\rm max},b_{\perp},P^z)$ from zero is a measure of the resulting truncation error. For the largest range of $z$ values we could realize numerically, $z_{\rm min}=-1.44$fm, $z_{\rm max}=1.44$fm, this error is still noticeable. This brute-force truncation of the FT leads to an oscillatory behavior of TMDWFs. This oscillation in $\tilde{\Phi}(x,b_{\perp},P^z)$ can be eliminated by an appropriate extrapolation for $\tilde{\Phi}(z,b_{\perp},P^z)$ as a function of $zP^z$ before Fourier transformation. While the signal-to-noise ratio of our data is not smooth enough, the brute-force Fourier transformation is adopted. 

%%%%%%%%%%%%%%%%%%%%%%%%%%%%%%%%%%%%%%
\begin{figure}
\centering
\includegraphics[width=0.62\textwidth]{qwf_co_b2}
\includegraphics[width=0.62\textwidth]{qwf_co_b4}
\includegraphics[width=0.62\textwidth]{qwf_co_i_b2}
\includegraphics[width=0.62\textwidth]{qwf_co_i_b4}
\caption{Examples for subtracted quasi TMDWFs in coordinate space. Here we take the cases of $P^z=24\pi/n_s$. According to our results, the subtracted quasi TMDWFs $\tilde{\Psi}^+(z,b_{\perp},P^z)$ are complex, thus the real and the imaginary part both need to be investigated. The upper four figures are for subtracted quasi TMDWFs $\tilde{\Psi}^+(z,b_{\perp},P^z)$ as a function of $\lambda=zP^z$ after the average over the two Dirac structures ($\gamma^t\gamma_5$ and $\gamma^z\gamma_5$) is taken. The upper left figure shows the real part of $\tilde{\Psi}^+(z,b_{\perp},P^z)$ with $b_{\perp}=2a$, while the upper right one is for $\tilde{\Psi}^+(z,b_{\perp},P^z)$ with $b_{\perp}=4a$. The lower two figures show the corresponding imaginary parts with $b_{\perp}=2a$ and $4a$.}
    \label{fig:qwf_co}
\end{figure}
%%%%%%%%%%%%%%%%%%%%%%%%%%%%%%%%%%%%%%


%%%%%%%%%%%%%%%%%%%%%%%%%%%%%%%%%%%%%%
\begin{figure}
\centering
\includegraphics[width=0.62\textwidth]{qwf_mom_b2}
\includegraphics[width=0.62\textwidth]{qwf_mom_b4}
\includegraphics[width=0.62\textwidth]{qwf_mom_i_b2}
\includegraphics[width=0.62\textwidth]{qwf_mom_i_b4}
\caption{Examples for subtracted quasi TMDWFs in momentum space with hadron momentum $P^z=24\pi/n_s$. The subtracted quasi TMDWFs in momentum space $\tilde{\Psi}^+(x,b_{\perp},P^z)$ are Fourier transformed from $\tilde{\Phi}^+(z,b_{\perp},P^z)$, which have real and imaginary part. Both have to be investigated. Shown in Eq. (\ref{eq:FT}), a brute-force FT is used to determine $\tilde{\Psi}^+(x,b_{\perp},P^z)$, where $z_{\rm min}=-1.44$~fm, $z_{\rm max}=1.44$~fm. The upper left figure shows the real part for $\tilde{\Psi}^+(x,b_{\perp},P^z)$ with $b_{\perp}=2a$ case, while the upper right one is for $\tilde{\Psi}^+(x,b_{\perp},P^z)$ with $b_{\perp}=4a$ case. Correspondingly, the lower two figures shows the imaginary part for $\tilde{\Psi}^+(x,b_{\perp},P^z)$ with $b_{\perp}=2a$ and $4a$.}
    \label{fig:qwf_momenta}
\end{figure}
%%%%%%%%%%%%%%%%%%%%%%%%%%%%%%%%%%%%%%


 


\subsection{Collins-Soper Kernel From Quasi TMDWFs}
\label{subsec:csresults}

The CS kernel governs the rapidity evolution and thus is independent of the momentum fraction of the involved parton.  But as indicated in Eq.(\ref{eq:matching}), the factorization formula  works only when $xP^z\gg \Lambda_{\rm QCD}$, and could be invalid  in the end-point regions $x\to0,1$. Power corrections are presumably  of the form $1/\left(xP^z\right)^2$ or $1/\left(\bar{x}P^z\right)^2$. Therefore,  the numerical CS kernel is fitted by a function of $x$, $P_1^z$ and $P_2^z$ and is written as $K(b_{\perp},\mu,x,P^z_1,P^z_2)$,
\begin{align}
&K(b_{\perp},\mu,x,P^z_1,P^z_2) =\nonumber\\
&\frac{1}{2\ln(P_1^z/P_2^z)} \left[\ln\frac{H^{+}(xP_2^z,\mu)\tilde{\Psi}^{+}(x,b_{\perp},\mu,P_1^z)}{H^{+}(xP_1^z,\mu)\tilde{\Psi}^{+}(x,b_{\perp},\mu,P_2^z)}\right.  \nonumber\\
	&+\left.\ln\frac{H^{-}(xP_2^z,\mu)\tilde{\Psi}^{-}(x,b_{\perp},\mu,P_1^z)}{H^{-}(xP_1^z,\mu)\tilde{\Psi}^{-}(x,b_{\perp},\mu,P_2^z)}\right].
\end{align}
Here $K(b_{\perp},\mu,x,P^z_1,P^z_2)$ are extracted from the perturbative matching kernels and quasi TMDWFs using 1-loop matching. They will have power corrections of teh form $\mathcal{O}\left(1/\left(xP^z\right)^2\right)$ and $\mathcal{O}\left(1/\left(\bar{x}P^z\right)^2\right)$. In order to extract the leading power contributions, we adopt  the following parametrization
\begin{align}
&K(b_{\perp},\mu,x,P^z_1,P^z_2) = K(b_{\perp},\mu) \nonumber\\
&\quad+A\left[\frac{1}{x^2(1-x)^2(P^z_1)^2}-\frac{1}{x^2(1-x)^2(P^z_2)^2}\right], \label{eq:parametrizationofCSkernel}
\end{align}
where $A$ is the coefficient accounting for  the leading higher power contributions, and can be determined through a joint fit of different lattice data in the regions not so close to $x=0, 1$. 

Fig.~\ref{fig:K_b} presents the physical CS kernel $K(b_{\perp},\mu)$ with $b_{\perp}=\{0.12,0.24,0.36,0.48,0.60\}$fm. By employing three cases of quasi TMDWFs with $P^z=\{8,10,12\}\times2\pi/n_s$, one can extract $K(b_{\perp},\mu,x,P^z_1,P^z_2)$ with $P_1^z/P_2^z=10/8$ and $12/8$, shown as the different colored bands.  Except in the end-point regions ($x<0.2$ or $x>0.8$), the lattice data is flat and reflects the leading power contribution, which conforms with expectations. Using the parametrization formula Eq.(\ref{eq:parametrizationofCSkernel}), the physical CS kernel $K(b_{\perp},\mu)$ can be determined by fitting the data, shown as the green band. As mentioned before, at large $b_{\perp}$, the quasi TMDWFs   show oscillations due to the  truncation of the Fourier transformation, which also affect the extracted $K(b_{\perp},\mu,x,P^z_1,P^z_2)$, as shown in the lower panel of Fig.\ref{fig:K_b}. This oscillation effect can be in principle removed once larger $z$ data becomes possible, or if one knows how to extrapolated the current data to the larger $z$ or $\lambda$ region. 


%%%%%%%%%%%%%%%%%%%%%%%%%%%%%%%%%%%%%%
\begin{figure}
\centering
    \includegraphics[width=0.62\textwidth]{K_b1}
    \includegraphics[width=0.62\textwidth]{K_b2}
    \includegraphics[width=0.62\textwidth]{K_b3}
    \includegraphics[width=0.62\textwidth]{K_b4}
    \includegraphics[width=0.62\textwidth]{K_b5}
\caption{The fit results of $K(b_{\perp},\mu,x,P^z_1,P^z_2)$ extracted from the quasi-WFs $\tilde{\Psi}^{\pm}$. The chosen momentum pairs $\{P_1^z,P^z_2\}$ are denoted by $P^z_1/P_2^z$ in the legend. The figures are corresponding to cases with $b_{\perp}=\{0.12,0.24,0.36,0.48,0.60\}$fm.} The horizontal shaded band shows the central value and uncertainty of $K(b_{\perp},\mu)$, as well as the fit range of $x$, as described in the text. In large $b_{\perp}$ area, the strong oscillation exists in the shaded area at both edges, where LaMET approach is invalid, is caused by the breakdown of the large momentum expansion.
    \label{fig:K_b}
\end{figure}
%%%%%%%%%%%%%%%%%%%%%%%%%%%%%%%%%%%%%%


Theoretically the physical CS kernel is purely real, however, there still exists a residual imaginary part at 1-loop matching. As discussed above, this imaginary term comes from the perturbative matching kernel. It is easily to prove that $H^{\pm}(zP_2^z,\mu)/H^{\pm}(zP_1^z,\mu)$ contains imaginary part while the lattice results for $\tilde{\Psi}^{\pm}(x,b_{\perp},\mu,P_1^z)/\tilde{\Psi}^{\pm}(x,b_{\perp},\mu,P_2^z)$ are nearly real, shown as \ref{fig:ratio_phi}. Therefore, we consider this imaginary part as systematic uncertainty from factorization theorem. This uncertainty can be expressed as

\begin{align}
&\sigma_{\mathrm{sys}}=\sqrt{K(b_{\perp},\mu)+\text{Im}^2\left[K^+(b_{\perp},\mu)\right]}-K(b_{\perp},\mu)
\label{eq:systematical_error}
\end{align}
where $\text{Im}\left[K^+(b_{\perp},\mu)\right]$ represents the numerical imaginary part of extracting $K(b_{\perp},\mu)$ only by $\tilde{\Psi}^+$:
\begin{align}
K^+(b_{\perp},\mu)=&\frac{1}{\ln (P^z_1/P^z_2)}\ln\frac{H^+(xP^z_2,\mu)\tilde \Psi^+(x, b_\perp, \mu, P^z_1)}{H^+(xP^z_1,\mu)\tilde \Psi^+(x, b_\perp, \mu, P^z_2)}.
\label{eq:K+}
\end{align}

It should be noticed that the perturbative matching kernel $H^+(zP^z,\mu)$ is the complex conjugate of $H^-(zP^z,\mu)$, that is the imaginary parts in these terms can be cancelled each other when we employ the average of $H^+(zP_2^z,\mu)/H^+(zP_1^z,\mu)$ and $H^-(zP_2^z,\mu)/H^-(zP_1^z,\mu)$. Therefore, as the final result, we adopt $K(b_{\perp},\mu)=\left[K^+(b_{\perp},\mu)+K^-(b_{\perp},\mu)\right]/2$ to reserve the real part, and regard the imaginary contributions as our systematic uncertainty.

%%%%%%%%%%%%%%%%%%%%%%%%%%%%%%%%%%%%%%
\begin{figure}
\centering
\includegraphics[width=0.62\textwidth]{imag_ratio_b3}
\caption{Example of numerical results for the imaginary part of the ratio for quasi TMDWFs at different $P^z$, $\frac{\tilde{\Psi}^+(x,b_{\perp},\mu,P_1^z)}{\tilde{\Psi}^+(x,b_{\perp},\mu,P_2^z)}$, with $b_{\perp}=3$a as a function of momentum fraction $x$. The imaginary parts of both cases are close to zero.}
    \label{fig:ratio_phi}
\end{figure}
%%%%%%%%%%%%%%%%%%%%%%%%%%%%%%%%%%%%%%

%%%%%%%%%%%%%%%%%%%%%%%%%%%%%%%%%%%%%%

\subsection{Results and Discussions}

One should notice that the Wilson loop renormalized quasi TMDWF on the lattice (Eq.(\ref{eq:WLRQuasi})) has a scale dependence on $a$. If one converts it to the $\overline{\rm MS}$ scheme through dividing it by $Z_{O}$ (Eq.(\ref{eq:ZO_RS_higher})), the scale $\mu$ is introduced. In principle, one should convert the Wilson loop renormalized quasi TMDWF to the $\overline{\rm MS}$ scheme since our factorization formula works there. However, since $Z_{O}$ has no dependence on momentum $P_{z}$, it cancels in the ratio of quasi TMDWFs, so does the scale dependence. So, one does not need to do the scheme and scale conversion of the quasi TMDWF during the extraction of CS kernel.

The extracted CS kernel from the combined fit of the ratios of quasi TMDWFs  with different momenta are shown as the red data points in Fig. \ref{fig:K_com}. {In this figure we exhibit two kinds of errors for $K(b_{\perp},\mu)$, in which the smaller ones denote statistical uncertainties while the larger ones include both statistical and systematical uncertainties}. In the small-$b_{\perp}$ region, systematical uncertainties are dominant due to the large power and the nonzero imaginary part. 

As a comparison, we also give the tree-level matching  result for the CS kernel.  With the leading order matching kernel $H(xP^z,\mu)=1+\mathcal{O}(\alpha_s)$, Eq. (\ref{eq:extractingCSkernel}) simplifies to the ratio of quasi TMDWFs $\tilde{\Phi}$ at $z=0$ with momentum $P_1^z/P^z_2$. The blue dots in Fig.~\ref{fig:K_com} denote the results obtained for tree-level matching, for which only statistical uncertainties are shown. 

%%%%%%%%%%%%%%%%%%%%%%%%%%%%%%%%%%%%%%
\begin{figure}
\centering
\includegraphics[width=0.62\textwidth]{K_com}
\includegraphics[width=0.62\textwidth]{K_pheno}
\caption{{The upper panel shows the comparison of our results $K(b_{\perp},\mu)$ and $K_0(b_{\perp},\mu)$ with the lattice calculations by SWZ~\cite{Shanahan:2021tst}, LPC~\cite{LatticeParton:2020uhz}, ETMC/PKU~\cite{Li:2021wvl} and SVZES\cite{Schlemmer:2021aij}, as well as the perturbative calculations up to 3-loop. $K(b_{\perp},\mu)$ denotes the CS kernel extracted through 1-loop matching, and whose uncertainties correspond to the statistical errors and the systematic ones from the non-zero imaginary part. $K_0(b_{\perp},\mu)$ denotes our tree-level results, only with statistical uncertainties. The lower panel shows the comparison of our result with phenomenological extractions: SV19~\cite{Scimemi:2019cmh}, Pavia19~\cite{Bacchetta:2019sam} and SIYY15~\cite{Sun:2014dqm} give phenomenological parameterizations of CS kernel fitted to data from} high energy collision processes like Drell-Yan.}
\label{fig:K_com}
\end{figure}
%%%%%%%%%%%%%%%%%%%%%%%%%%%%%%%%%%%%%%

We compare our results with the ones from perturbative calculations, phenomenological extractions as well as the lattice results determined by other collaborations. 

The black solid and dashed lines in the upper panel of Fig.~\ref{fig:K_com} indicate the perturbative results  up to 3-loops, with a running coupling constant $\alpha_s(\mu=1/b_{\perp})$. The perturbative calculations work well in small $b_{\perp}$ region ($b_{\perp}\ll1/\Lambda_{\mathrm{QCD}}$), while will diverge with $b_{\perp}$ increasing. In contrast, the lattice calculation will give accurate predictions in the nonperturvative region, while due to the power corrections, it 
might suffer large systematic uncertainties in small $b_{\perp}$ region.

Similar to this work, the results of LPC~\cite{LatticeParton:2020uhz} and ETMC/PKU~\cite{Li:2021wvl} are also extracted from quasi TMDWFs through a tree-level matching. Adopting the one-loop matching formula, as well as considering the operator mixing effects will help ones to reduce the systematic uncertainties and obtain more precise results. In addition, considering the different directions of gauge link will help us to eliminate the contributions from unphysical imaginary part, and then improve the accuracy of our results.

In another way, the SWZ~\cite{Shanahan:2021tst} and SVZES\cite{Schlemmer:2021aij} results are obtained from quasi-TMDPDFs. Compared with the complicated nucleon correlation functions, the meson ones are much easier to obtain better signals.  Besides, the wave functions of meson are nearly symmetric in $x$-space, it thereby more convenient to parametrize the oscillation effects and obtain the physical results in large $P^z$ limit like Fig.(\ref{fig:K_b}). In addition, the light meson is more easily to reach a larger boosted factor, one can see from the small $b_{\perp}$ region, the results from quasi TMDWFs fit well with the perturbative calculations than the ones from quasi TMDPDFs.

The lower panel of Fig.(\ref{fig:K_com}) shows the comparison with phenomenological results. SV19~\cite{Scimemi:2019cmh} and SIYY15~\cite{Sun:2014dqm} use a parameterization with perturbative and nonperturbative parts. However Pavia19~\cite{Bacchetta:2019sam} obtained their result with the factorization of TMDPDFs, obtaining the CS kernel from the rapidity derivative. In addition they fit parameters from the Drell-Yan data to obtain their phenomenological CS kernel. The results from different methods exhibit obviously inconsistency in the nonperturbative region. Our result shows a better consistency with SV19.
